\documentclass[10pt,twocolumn]{article}
\usepackage[margin=0.6in]{geometry}
\usepackage{float}
\usepackage{setspace}
\usepackage[compact,small]{titlesec}
\usepackage{enumitem}
\setstretch{1.75}
\setlist{nosep,leftmargin=*}
\setlength{\columnsep}{0.25in}
\setlength{\textfloatsep}{6pt}
\setlength{\abovecaptionskip}{2pt}
\titlespacing*{\section}{0pt}{4pt}{0pt}
\titlespacing*{\subsection}{0pt}{2pt}{0pt}
\renewcommand{\arraystretch}{1.3}

\title{\vspace{-0.4in}CS 431/531 Homework 2}
\author{Eleanor Stenberg}
\date{\vspace{-0.3in}}

\begin{document}
\maketitle

\section{Project overview}
In this project, we are computing a prefix sum, which replaces each array element with the total of itself and everything before it. The serial version is keeping a continual total ($N-1$ additions), but each answer depends on the last one, so it can't be split up directly. We build up partial sums in a tree, then pass them back down to fill in the missing values. 


\section{Parallel Versions}
For this project there are two different two parallel versions. Each utilize pragma omp parallel for in their steps to split each of the steps across threads. There is an O(N log N) version and 2(N-1) version.

The O(N log N) version has every element adding the value d spots trailing it, with d doubling on each step. The 2(N-1) version is completed in two parts. It starts by adding pairs of numbers, then onto pairs of pairs, and repeating until the very last spot holds the total. It then goes back down and goes to fill in any spots that were skipped.


\section{Project Setup}
The project was compiled on my MacBook Pro, using \texttt{gcc-16 -O3 -fopenmp} and ran on my MacBook Pro. I used 125 MB per array, and for each round, I utilized 1,2,4,6,8, and 10 threads to run it. Every run was successful in passing, showing that both parallel versions gave the same answers as the serial. 

\section{Results}

\begin{table}[H]\centering\small
\begin{tabular}{l|c|c}
Version & Additions (in millions) & vs.\ serial \\ \hline
Serial     & 34  & $1\times$  \\
O(N log N) & 805 & $24\times$ \\
2(N-1)     & 67  & $2\times$  \\
\end{tabular}
\end{table}

\begin{table}[H]\centering\small
\begin{tabular}{c|c|c|c}
Threads & Serial & O(N log N) & 2(N-1) \\ \hline
1  & 0.0388 & 0.1172 & 0.0485 \\
2  & 0.0387 & 0.0781 & 0.0364 \\
4  & 0.0380 & 0.0754 & 0.0314 \\
6  & 0.0394 & 0.0794 & 0.0297 \\
8  & 0.0369 & 0.0820 & 0.0322 \\
10 & 0.0393 & 0.0805 & 0.0294 \\
\end{tabular}
\end{table}

\begin{table}[H]\centering\small
\begin{tabular}{c|c|c}
Threads & O(N log N) & 2(N-1) \\ \hline
1  & 0.33 & 0.80 \\
2  & 0.50 & 1.06 \\
4  & 0.50 & 1.21 \\
6  & 0.50 & 1.33 \\
8  & 0.45 & 1.14 \\
10 & 0.49 & 1.34 \\
\end{tabular}
\end{table}

\begin{table}[H]\centering\small
\begin{tabular}{c|c|c|c}
Threads & O(N log N) & 2(N-1) & Perfect \\ \hline
1  & 1.00 & 1.00 & 1  \\
2  & 1.50 & 1.33 & 2  \\
4  & 1.55 & 1.54 & 4  \\
6  & 1.48 & 1.63 & 6  \\
8  & 1.43 & 1.50 & 8  \\
10 & 1.46 & 1.65 & 10 \\
\end{tabular}
\end{table}

\section{Results}
Through all versions, the 2(N-1) version proved to be the fastest, with it's best time being 0.0294 seconds using 10 threads. This ran roughly 1.34 times faster than the serial version. The serial ran faster than the O(N log N) version. 

Neither of the parallel versions got noticeably different in speed after 4 threads. 10 threads proved the fastest speed.

\section{Conclusion}
In conclusion, a paralell algorithm is only helpful if it does not have that much extra work that it is adding. 

\end{document}
